\subsection{Quantum Fourier transform}\glsadd{cc:FT:QFT}

Recall the tensor network decomposition of the Fourier transform developed in \secref{sec:fourierTransform}.

When $\vectorat{\shortcatvariables}$ is a quantum state, the Fourier transform can be understood as a unitary evolving the state.
Since both the controlled $U$ tensor and the Hadamard $H$ are unitary, each of them can be implemented by a sequence of gates.
For sparse implementation one can exploit that $U$ is a tensor product of matrices (i.e. single qubit gates).

% Gate representation
We introduce the gates
\begin{align*}
    R^{\catenumerator}[\ccatvariable{\insymbol},\ccatvariable{\outsymbol}] =
    \coloredmatrixof{
        1 & 0 \\
        0 & \expof{2\pi i \frac{1}{2^\catenumerator}}
    }{\ccatvariable{\insymbol},\ccatvariable{\outsymbol}} \, .
\end{align*}
The unitary $U^{\catorder}$ in the tensor network decomposition of the fourier transform (see \secref{sec:ft:tn}) is then the tensor product
\begin{align*}
    U^{\catorder}[\ccatvariableof{\insymbol}{[\catorder]},\ccatvariableof{\outsymbol}{[\catorder]}]
    = \bigotimes_{\catenumeratorin} R^{\catenumerator+2}[\ccatvariableof{\insymbol}{\catenumerator},\ccatvariableof{\outsymbol}{\catenumerator}] \, .
\end{align*}
Controlling $U^{\catorder}$ by $\catvariableof{\catorder}$ is equivalent to controlling each $R^{\catenumerator}$ gate.


% Gate complexity
To realize the QFT we need $\catorder$ single-qubit Hadamard gates and $\sum_{\catenumeratorin} \catenumerator = \frac{\catorder\cdot(\catorder-1)}{2}$ controlled diagonal single-qubit gates (to realize the $U^{\catenumerator}$ transforms).
The gate complexity is therefore $\mathcal{O}(\catorder^2)$.

\subsubsection{Walsh-Hadamard transform}\glsadd{qc:WH}

In case of $\catdimof{\catenumerator}=2$, the Fourier transform is known to be the Walsh-Hadamard tranform.
The transform is performed by applying Hadamard gates on each variable.
This is equivalent with contraction of the tensor
\begin{align*}
    \hgate^{\catorder}\left[\catvariableof{[\catorder],\insymbol},\catvariableof{[\catorder],\outsymbol}\right] =
    \bigotimes_{\catenumeratorin} \hgateat{\ccatvariableof{\catenumerator}{\insymbol},\ccatvariableof{\catenumerator}{\outsymbol}} \, ,
\end{align*}
which coordinates are
\begin{align*}
    \hgate^{\catorder}\left[\indexedcatvariableof{[\catorder],\insymbol},\indexedcatvariableof{[\catorder],\outsymbol}\right] =
    \sqrt{\frac{1}{2^{\catorder}}} \left(-1\right)^{\sum_{\catenumeratorin}\catindexof{\catenumerator,\insymbol} \cdot \catindexof{\catenumerator,\outsymbol}} \, .
\end{align*}

We have
\begin{align*}
    \contractionof{\onehotmapofat{0}{\ccatvariable{\insymbol}},
        \hgateat{\ccatvariable{\insymbol},\ccatvariable{\outsymbol}}}{\ccatvariable{\outsymbol}}
    &= \sqrt{\frac{1}{2}} \onesat{\ccatvariable{\outsymbol}} \\
    \contractionof{\onehotmapofat{1}{\ccatvariable{\insymbol}},
        \hgateat{\ccatvariable{\insymbol},\ccatvariable{\outsymbol}}}{\ccatvariable{\outsymbol}}
    &= \sqrt{\frac{1}{2}} \coloredmatrixof{1 \\ -1}{\ccatvariable{\outsymbol}}
\end{align*}
and conversely
\begin{align*}
    \contractionof{\sqrt{\frac{1}{2}} \onesat{\ccatvariable{\insymbol}},
        \hgateat{\ccatvariable{\insymbol},\ccatvariable{\outsymbol}}}{\ccatvariable{\outsymbol}}
    &= \onehotmapofat{0}{\ccatvariable{\outsymbol}} \\
    \contractionof{\sqrt{\frac{1}{2}} \coloredmatrixof{1 \\ -1}{\ccatvariable{\insymbol}},
        \hgateat{\ccatvariable{\insymbol},\ccatvariable{\outsymbol}}}{\ccatvariable{\outsymbol}}
    &= \onehotmapofat{1}{\ccatvariable{\outsymbol}} \, .
\end{align*}
If and only if the function is constant, then the transformed state is the ground state.

Let us present two known algorithms as the measurement of states prepared by \computationCircuits{}.

\begin{example}[Deutsch-Jozsa Test]
    We measure the probability of the ground state of the sign encoding after a Walsh-Hadamard transform.
    The probability of measuring the ground state after a Walsh-Hadamard transform of the distributed qubits is then
    \begin{align*}
        \frac{1}{2^{\catorder}} \cdot \contraction{\absof{\contractionof{\signqstateofat{\exformula}{\shortcatvariables,\headvariable}}{\headvariable}}^2}
        = \frac{1}{2^{2\cdot\catorder}} \left(\contraction{\exformula}-\contraction{\lnot\exformula}\right)^2 \, .
    \end{align*}
    Thus, if and only if the probability of the ground state is $1$, we have that the formula is constant, that is $\exformula\in\{\top,\bot\}$.
    If and only if the probability of the ground state is $0$, we have that the formula is balanced, that is $\contraction{\exformula}=\contraction{\lnot\exformula}$.
    When its known, that the formula is either constant or balanced, we can decide which one it is by measuring the probability of the ground state.
    This is the Deutsch-Josza algorithm \cite{deutsch_rapid_1992}.

    %%
    Note that when allowing for an exponentially small error the decision can be efficiently performed by a classical algorithm.
    When drawing $n$ samples and deciding on $\meanparam=0.5$ if and only if in the samples a $0$ and a $1$ is observed, we have an error probability of $0$ when $\meanparam\in\{0,1\}$ and of
    \begin{align*}
        \probof{E|\meanparam=0.5} = 2 \cdot \frac{1}{2^{n}} \, .
    \end{align*}
    The quantum advantage is thus bound to the demand of an error-free decision, for which a generic classical algorithm would iterate through all states and is therefore exponentially slower than the quantum alternative.

    \begin{figure}
        \begin{center}

            \begin{tikzpicture}[scale=0.35,thick]

                \draw (0,0) rectangle (2,7.5);
                \node[anchor=center] (text) at (1,3.75) {\corelabelsize $\signqstateof{\exformula}$};

                \draw (2,6.5) -- (4,6.5) node[midway,above] {\colorlabelsize $\headvariable$};
                \drawqcmeasuresymbol{5}{6.5}
                \draw (6,6.5) -- (8,6.5); % node[midway,above] {\colorlabelsize $\headvariable$};
                \drawvariabledot{8}{6.5}

                \draw (2,4) -- (4,4) node[midway,above] {\colorlabelsize $\catvariableof{0}$};
                \draw (4,3) rectangle (6,5);
                \node[anchor=center] at (5,4) {$\hgate$};
                \draw (6,4) -- (7,4);
                \drawqcmeasuresymbol{8}{4}
                \draw (9,4) -- (10,4);
                \draw (10,3) rectangle (12,5);
                \node[anchor=center] at (11,4) {$\fbasis$};

                \node[anchor=center] (text) at (3,3.25) {$\vdots$};
                \draw (2,1) -- (4,1) node[midway,above] {\colorlabelsize $\catvariableof{\catorder\shortminus 1}$};
                \draw (4,0) rectangle (6,2);
                \node[anchor=center] at (5,1) {$\hgate$};
                \draw (6,1) -- (7,1);
                \drawqcmeasuresymbol{8}{1}
                \draw (9,1) -- (10,1);
                \draw (10,0) rectangle (12,2);
                \node[anchor=center] at (11,1) {$\fbasis$};

                \node[anchor=center] (text) at (14.5,3.75) {${=}\,\, \frac{1}{2^{2\cdot\catorder}}$};

                \begin{scope}[shift={(18,1.5)}]

                    % Squared
                    \draw (-1,6) to[bend right=15] (-1,-2);
                    \draw (7,6) to[bend right=-15] (7,-2);
                    \drawtextnode{7.5}{6.5}{$2$}

                    \draw (1.5,3) rectangle (4.5,6);
                    \drawtextnode{3}{4.5}{
                        $\begin{bmatrix}
                             -1 \\
                             1
                        \end{bmatrix}$
                    }

                    \draw[->-] (3,1) -- (3,3) node[midway,right] {\colorlabelsize $\headvariable$};
                    \draw (0,-1) rectangle (6,1);
                    \drawtextnode{3}{0}{\corelabelsize $\bencodingof{\exformula}$}
                    \draw[->-] (1,-2) -- (1,-1) node[midway,left] {\colorlabelsize $\catvariableof{0}$};
                    \drawvariabledot{1}{-2}
                    \drawtextnode{3}{-1.5}{$\cdots$}
                    \draw[->-] (5,-2) -- (5,-1) node[midway,right] {\colorlabelsize $\catvariableof{\catorder\shortminus1}$};
                    \drawvariabledot{5}{-2}
                \end{scope}

            \end{tikzpicture}
        \end{center}
        \caption{Measurement setup for the Deutsch-Jozsa algorithm on the formula $\exformula$.
        }\label{fig:deutschJozsa}
    \end{figure}

\end{example}


% Generalization
\todo{Another generalizations would be period detection.}
Deutsch-Jozsa can be generalized to estimating the moments of a distribution given its q-sample.
Note, that we have to demand a q-sample (i.e. vanishing phase tensor), otherwise the measurement is affected by the phase tensor and the Deutsch-Jozsa test is effectively done wrt another function.
If for example the coordinates of the phase tensor are multiples of $\pi$, and the phases are thus $+1/-1$, then the test is done wrt the function $\exformula\oplus(\frac{1}{\pi}\phasecore)$ instead of $\exformula$.

\begin{example}[Inversion Test]
    Concatenating the \computationCircuits{} to two formulas $\exformula$ and $\secexformula$ prepares the state
    \begin{align*}
        \basqstateofat{\exformula\oplus\secexformula}{\shortcatvariables,\headvariable}
    \end{align*}
    The probability of measuring the ground state after a Walsh-Hadamard transform of the distributed qubits (see \figref{fig:inversionTest}) is then
    \begin{align*}
        \frac{1}{2^{\catorder}} \absof{\contractionof{\basqstateofat{\exformula\oplus\secexformula}{\shortcatvariables,\headvariable}}{\headvariable}}^2
        &= \frac{1}{2^{2\cdot\catorder}}
        \begin{bmatrix}
            \contraction{\exformula\oplus\secexformula}
            \contraction{\lnot(\exformula\oplus\secexformula)}
        \end{bmatrix} \\
        &= \frac{1}{2^{2\cdot\catorder}}
        \begin{bmatrix}
            \#{\left\{\shortcatindices\in\atomstates \wcols \exformulaat{\indexedshortcatvariables}=\secexformula\left[\indexedshortcatvariables\right]\right\}} \\
            \#{\left\{\shortcatindices\in\atomstates \wcols \exformulaat{\indexedshortcatvariables}\neq\secexformula\left[\indexedshortcatvariables\right]\right\}}
        \end{bmatrix} \, .
    \end{align*}
    The second equation holds since $\lnot(\exformula\oplus\secexformula)$ is equivalent with the biconditional $\exformula\Leftrightarrow\secexformula$.

    % Inversion test
    This is in more generality the inversion test of measuring the overlap of $\basqstateof{\exformula}$ with $\basqstateof{\secexformula}$ (see \cite{schuld_machine_2021}).
    Here we used that the \computationCircuit{} is self-adjoint and understand the concatenation as the unitaries preparing $\basqstateof{\exformula}$ and the adjoint (together with the Walsh-Hadamard transform) of $\basqstateof{\secexformula}$.

    \begin{figure}
        \begin{center}

            \begin{tikzpicture}[scale=0.35,thick]

                \draw (0,0) rectangle (2,7.5);
                \node[anchor=center] (text) at (1,3.75) {\corelabelsize $\basqstateof{\exformula\oplus\secexformula}$};

                \draw (2,6.5) -- (4,6.5) node[midway,above] {\colorlabelsize $\headvariable$};
                \drawqcmeasuresymbol{5}{6.5}
                \draw (6,6.5) -- (8,6.5) node[midway,above] {\colorlabelsize $\headvariable$};

                \draw (2,4) -- (4,4) node[midway,above] {\colorlabelsize $\catvariableof{0}$};
                \draw (4,3) rectangle (6,5);
                \node[anchor=center] at (5,4) {$\hgate$};
                \draw (6,4) -- (7,4);
                \drawqcmeasuresymbol{8}{4}
                \draw (9,4) -- (10,4);
                \draw (10,3) rectangle (12,5);
                \node[anchor=center] at (11,4) {$\fbasis$};

                \node[anchor=center] (text) at (3,3.25) {$\vdots$};
                \draw (2,1) -- (4,1) node[midway,above] {\colorlabelsize $\catvariableof{\catorder\shortminus 1}$};
                \draw (4,0) rectangle (6,2);
                \node[anchor=center] at (5,1) {$\hgate$};
                \draw (6,1) -- (7,1);
                \drawqcmeasuresymbol{8}{1}
                \draw (9,1) -- (10,1);
                \draw (10,0) rectangle (12,2);
                \node[anchor=center] at (11,1) {$\fbasis$};

                \node[anchor=center] (text) at (14.5,3.75) {${=}\,\, \frac{1}{2^{2\cdot\catorder}}$};

                \begin{scope}[shift={(18,3.75)}]

                    \draw (3,-0.5) ellipse (4.5cm and 3cm);
                    \node[anchor=center] at (5,3.25) {$\absof{\cdot}^2$};

                    \draw[->-] (3,1) -- (3,3) node[midway,right] {\colorlabelsize $\headvariable$};
                    \draw (0,-1) rectangle (6,1);
                    \drawtextnode{3}{0}{\corelabelsize $\bencodingof{\exformula\oplus\secexformula}$}
                    \draw[->-] (1,-2) -- (1,-1) node[midway,left] {\colorlabelsize $\catvariableof{0}$};
                    \drawvariabledot{1}{-2}
                    \drawtextnode{3}{-1.5}{$\cdots$}
                    \draw[->-] (5,-2) -- (5,-1) node[midway,right] {\colorlabelsize $\catvariableof{\catorder\shortminus1}$};
                    \drawvariabledot{5}{-2}
                \end{scope}

            \end{tikzpicture}
        \end{center}
        \caption{Measurement setup for the inversion test on the formulas $\exformula$ and $\secexformula$.
        We measure the probability of the statistic qubit $\headvariable$, when the distributed qubits $\shortcatvariables$ are in the ground state after a Walsh-Hadamard transform.
        This is equal to a scaled and squared contraction of the basis encoding to the formula $\exformula\oplus\secexformula$.
        }\label{fig:inversionTest}
    \end{figure}
\end{example}


%\subsection{Further Overlap-measuring Circuits}
Quantum Circuits such as the SWAP test and the Hadamard test (\cite[Section 3.6.1]{schuld_machine_2021}) can be used to measure overlaps of quantum states, which are the squared absolutes of contractions of two state tensors.
\begin{itemize}
    \item When we have preparation schemes for two tensors, we can control them with a common ancilla qubit and measure their contraction by a Hadamard test (alternatively, using the SWAP test and state preparation in two registers).
    \item Can we extend these schemes to contractions of more general tensor networks?
\end{itemize}

%\subsection{Discussion}
While the preparation of contracted basis encodings is done efficiently, accessing the prepared state by measurements is a bottleneck.
This is a typical \emph{quantum state tomography} problem \cite{roth_semi-device-dependent_2023}.
Sparsity of the contracted state can be exploited to probabilistically guarantee precise estimates \cite{gross_quantum_2010}.
However, when deciding whether $\exformula$ is satisfiable, based on the prepared contraction $\sqrt{\normalizationof{\bencodingofat{\exformula}{\headvariable,\shortcatvariables}}{\headvariable}}$, we have an exponentially small error tolerance.
Moreover, the applied sample mean is the UMVUE (Uniformly Minimum Variance Unbiased Estimator, see \cite{casella_statistical_2001}), and can thus not be improved by other unbiased estimators.




\subsubsection{Period detection}

QFT is a typical component in Quantum algorithms: QFT can be done efficiently with $\mathcal{0}(\catorder^2)$ controlled single-qubit gates.
If the Fourier transform is a basis vector, a single measurement is sufficient to determine it.
We now investigate how this can be used in the detection of periods.

%\subsubsection{Detecting periods}

Based on the support of the Fourier transform, one can derive the periods of a function.
Periods are indices $(\thirdcatvariableof{0},\ldots,\thirdcatvariableof{\catorder-1})$ such that
\begin{align*}
    \hypercoreat{\shortcatvariables=\shortcatindices}
    = \hypercoreat{\shortcatvariables=\shortcatindices\oplus\thirdcatindexof{[\catorder]}}
\end{align*}
where $\oplus$ is the position-wise sum $\mathrm{mod}\,\,\catdimof{\catenumerator}$.


Each supported index of the Fourier transform corresponds with periods in the function.
The periods of the functions are thus at least the intersection of those (when assuming more structure one might have more periods).

%% WH-case
In case of the Walsh-Hadamard transform (i.e. $\catdimof{\catenumerator}=2$), we can easily determine the periods corresponding to each basis vector.
The inverse Fourier transform of $\onehotmapofat{\selindexof{[\catorder]}}{\selvariableof{[\catorder]}}$ is the tensor
\begin{align*}
    \left(\bigotimes_{\catenumeratorin \wcols \selindexof{\catenumerator}=0} \hplusbasat{\catvariableof{\catenumerator}} \right)
    \otimes
    \left(\bigotimes_{\catenumeratorin \wcols \selindexof{\catenumerator}=1} \hminusbasat{\catvariableof{\catenumerator}} \right) \, .
\end{align*}

For each $\hplusbasat{\catvariableof{\catenumerator}}$ we have a period generator of $\onehotmapof{\catenumerator}$.
For each pair of $\hminusbasat{\catvariableof{\catenumerator}}\otimes \hminusbasat{\catvariableof{\seccatenumerator}}$ we have a period generator of $\onehotmapof{\catenumerator}+\onehotmapof{\seccatenumerator}$.

One example is the Deutsch-Jozsa test, where the transformed is the $\onehotmapofat{0}{\selvariableof{[\catorder]}}$ state, if and only if the function is constant.
In this case, any vector is a period, since the group itself is generated by the period generators.


%\subsubsection{Generic Hidden Subgroup Problems}

% Generic hidden subgroup problems
%In general \cite{nielsen_quantum_2010} approached by generic Fourier transforms of the distributed variables of a basis encoding state to a function.